\documentclass[11pt,a4paper]{article}

% ─── Packages ────────────────────────────────────────────────────────────────
\usepackage[T1]{fontenc}
\usepackage[utf8]{inputenc}
\usepackage{lmodern}
\usepackage[margin=2.5cm]{geometry}
\usepackage{amsmath,amssymb,amsthm}
\usepackage{bm}
\usepackage{booktabs}
\usepackage{array}
\usepackage{xcolor}
\usepackage{hyperref}
\usepackage{framed}

% ─── Colours ─────────────────────────────────────────────────────────────────
\definecolor{biogreen}{RGB}{30,120,70}
\definecolor{mathblue}{RGB}{10,60,140}
\definecolor{boxframe}{RGB}{100,130,200}
\definecolor{shadecolor}{RGB}{245,248,255}  % used by framed's snugshade

% ─── Hyperref setup ──────────────────────────────────────────────────────────
\hypersetup{
  colorlinks=true,
  linkcolor=mathblue,
  urlcolor=mathblue,
  citecolor=mathblue
}

% ─── Theorem environments ────────────────────────────────────────────────────
\theoremstyle{definition}
\newtheorem{definition}{Definition}[section]
\newtheorem{remark}{Remark}[section]
\newtheorem{example}{Example}[section]

\theoremstyle{plain}
\newtheorem{proposition}{Proposition}[section]

% ─── Notation shorthands ─────────────────────────────────────────────────────
\newcommand{\sig}[2]{S_{#1}(#2)}
\newcommand{\dS}[2]{\Delta S_{#1}(#2)}
\newcommand{\Nb}[2]{\mathcal{N}_{#1}(#2)}
\newcommand{\JE}[2]{\mathcal{J}_{#1}(#2)}
\newcommand{\FJ}[2]{\mathcal{F}_{#1}(#2)}
\newcommand{\expo}[2]{E_{#1}(#2)}
\newcommand{\RD}{\mathrm{RD}}
\newcommand{\RR}{\mathrm{RR}}
\newcommand{\pE}{p_{\mathrm{exp}}}
\newcommand{\pU}{p_{\mathrm{unexp}}}
\newcommand{\1}[1]{\mathbf{1}\!\left[#1\right]}

% ─────────────────────────────────────────────────────────────────────────────
\begin{document}

\begin{center}
  {\LARGE\bfseries\color{mathblue}
    Spatiotemporal Signal Propagation\\[4pt]
    in Live-Cell Imaging Data}\\[10pt]
  {\large A Mathematical Formalization}\\[6pt]
  {\normalsize Systems Biology, University of Warsaw}\\[4pt]
  {\small\color{gray}
    Companion to \texttt{spatiotemporal\_signal\_propagation.py}
    and \texttt{compare\_spatiotemporal\_behavior.py}}
\end{center}

\vspace{0.5em}
\noindent\rule{\linewidth}{0.8pt}
\vspace{1em}

\tableofcontents
\bigskip
\noindent\rule{\linewidth}{0.4pt}
\bigskip

% ─────────────────────────────────────────────────────────────────────────────
\section{Biological Setting and Motivation}
% ─────────────────────────────────────────────────────────────────────────────

The data originate from live-cell fluorescence microscopy experiments on
\textbf{MCF10A} cells (a non-tumorigenic human mammary epithelial line). Each
cell carries two fluorescent kinase-activity reporters:

\begin{itemize}
  \item \textbf{ErkKTR} --- an ERK kinase translocation reporter whose
        cytoplasmic-to-nuclear fluorescence ratio (\texttt{ERKKTR\_ratio})
        is a real-time proxy for ERK activity.
  \item \textbf{FoxO3A} --- a nuclear/cytoplasmic ratio reporter for the
        PI3K$\to$AKT$\to$FoxO3A signalling axis.
\end{itemize}

Cells were growth-factor-deprived for 48~h prior to acquisition, seeded in
96-well plates ($\sim 25\,000$ cells/well), and imaged every
$\Delta t = 5$~min for a total of $T_{\max} = 24$~h (up to $288$ frames).
The dataset spans six replicate experiments and $28$ imaging sites assigned to
different \emph{PI3K-pathway genetic backgrounds}: wild-type (WT),
\textbf{AKT1}~E17K, \textbf{PIK3CA}~E545K, \textbf{PIK3CA}~H1047R, and
\textbf{PTEN} deletion.

\bigskip
\begin{snugshade}
\noindent\textbf{Core biological question.}
\emph{When a cell shows a sudden, large increase (``jump'') in ERK or FoxO3A
activity, are its spatial neighbours more likely to show their own large
activity increase shortly afterwards?}
If yes, this suggests local cell-to-cell signalling or coordinated behaviour,
and we want to quantify how strongly this effect depends on the genetic
background.
\end{snugshade}

% ─────────────────────────────────────────────────────────────────────────────
\section{Data Model: Cells, Tracks, and Nodes}
% ─────────────────────────────────────────────────────────────────────────────

\subsection{Primary observables}

Let $\mathcal{C}$ be the set of all \emph{cell tracks} in one
experiment-site block. A track $k \in \mathcal{C}$ is a sequence of
time-resolved observations of the \emph{same physical cell}.
Let $T_k \subseteq \{1, 2, \ldots, T_{\max}\}$ be the set of frames during
which track $k$ is observed.

\begin{definition}[Node]\label{def:node}
A \emph{node} is a pair $v = (k,t)$ with $k \in \mathcal{C}$ and
$t \in T_k$. The set of all nodes is
\[
  V \;=\; \bigl\{(k,t) \;\big|\; k \in \mathcal{C},\; t \in T_k\bigr\}.
\]
Each node $(k,t)$ carries:
\begin{align*}
  \bm{x}(k,t) &= \bigl(x_{k,t},\, y_{k,t}\bigr) \in \mathbb{R}^2
    && \text{nuclear centroid (image pixels),}\\
  \sig{k}{t}  &\in \mathbb{R}_{>0}
    && \text{fluorescence ratio (signal channel),}\\
  A(k,t)      &\in \mathbb{R}_{>0}
    && \text{nuclear area (pixels).}
\end{align*}
\end{definition}

\begin{remark}
Physical time in hours is $t_{\mathrm{h}} = t \cdot \Delta t / 60$
with $\Delta t = 5$~min.
\end{remark}

% ─────────────────────────────────────────────────────────────────────────────
\section{Graph Construction}
% ─────────────────────────────────────────────────────────────────────────────

The scripts build a \emph{spatiotemporal graph}
$G = (V,\, E_{\mathrm{sp}} \cup E_{\mathrm{tm}})$
with two disjoint families of edges.

\subsection{Spatial edges}

\begin{definition}[Spatial neighbourhood]\label{def:spatial-nb}
For node $(k,t)$ and radius $r > 0$, the \emph{spatial neighbourhood} is
\[
  \Nb{k}{t} \;=\;
  \bigl\{(k',t) \;\big|\; k' \neq k,\;
    \|\bm{x}(k,t) - \bm{x}(k',t)\|_2 \leq r\bigr\}.
\]
A \emph{spatial edge} exists between $(k,t)$ and $(k',t)$ whenever
$(k',t) \in \Nb{k}{t}$, i.e.\ both cells are within distance $r$ of each
other \emph{at the same frame $t$}. The full set is
\[
  E_{\mathrm{sp}} \;=\;
  \bigl\{\{(k,t),(k',t)\} \;\big|\; (k',t) \in \Nb{k}{t}\bigr\}.
\]
\end{definition}

The default radius is $r = 60$ image-coordinate units. Spatial edges are
\emph{undirected} and \emph{frame-local}: a cell pair is connected only at
the frame(s) at which they are close.

\begin{remark}[Efficient computation via $k$-d trees]
At each frame $t$, all nuclear centroids are inserted into a
\emph{$k$-d tree} (a balanced binary space-partition structure).
Querying all pairs within radius $r$ costs
$\mathcal{O}(|V_t|\log|V_t| + |E_{\mathrm{sp},t}|)$ per frame,
where $V_t = \{(k,t) \mid t \in T_k\}$.
The $k$-d tree approach is therefore far more efficient than naively
checking all $\binom{|V_t|}{2}$ pairs.
\end{remark}

\subsection{Temporal edges}

\begin{definition}[Temporal edge]\label{def:temporal-edge}
For each track $k$ and each pair of \emph{consecutive observed frames}
$t, t' \in T_k$ with $t < t'$ and no intervening observation in $T_k$,
there is a \emph{directed temporal edge} $(k,t) \to (k,t')$. The set is
\[
  E_{\mathrm{tm}} \;=\;
  \bigl\{(k,t) \to (k,t') \;\big|\; k \in \mathcal{C},\;
    t' = \min_{s \in T_k,\; s > t} s\bigr\}.
\]
\end{definition}

Each temporal edge records the frame gap $\delta = t' - t$ (nominally
$\delta = 1$, but tracking gaps can exceed one frame) and the physical
time gap $\delta \cdot \Delta t$ in minutes.

\begin{remark}[Graph structure]
The temporal edges alone form a collection of directed paths --- one per
track --- corresponding to the trajectories of individual cells over time.
Adding spatial edges introduces \emph{cross-track} connections at each
frame, converting the path collection into a richer spatiotemporal graph.
\end{remark}

% ─────────────────────────────────────────────────────────────────────────────
\section{Signal Dynamics and Jump Detection}
% ─────────────────────────────────────────────────────────────────────────────

\subsection{First difference (signal delta)}

Along each temporal edge $(k,t) \to (k,t')$ in $E_{\mathrm{tm}}$, define the
\emph{signal increment}
\[
  \dS{k}{t'} \;=\; \sig{k}{t'} - \sig{k}{t}.
\]
When $(k,t)$ is the first observed frame of track $k$ (i.e.\ $(k,t)$ has no
incoming temporal edge), $\dS{k}{t}$ is undefined (\texttt{NaN}).

\subsection{Data-driven jump threshold}

\begin{definition}[Jump threshold]\label{def:threshold}
Let
\[
  \mathcal{D}^+ \;=\;
  \bigl\{\dS{k}{t} \;\big|\; (k,t) \in V,\; \dS{k}{t} > 0\bigr\}
\]
be the multiset of all \emph{positive} signal increments in the block.
For a chosen quantile $q \in (0,1)$, the \emph{data-driven jump threshold}
is
\[
  \theta_q \;=\; Q_q\!\bigl(\mathcal{D}^+\bigr),
\]
where $Q_q$ is the $q$-th empirical quantile.
The default is $q = 0.90$, so $\theta_{0.90}$ equals the 90th percentile of
all upward steps in the block.
A user-supplied fixed threshold overrides this per-block estimate.
\end{definition}

\subsection{Jump event indicator}

\begin{definition}[Jump event]\label{def:jump}
For node $(k,t)$, the \emph{jump event indicator} is
\[
  \JE{k}{t} \;=\; \1{\dS{k}{t} \geq \theta},
\]
using the convention $\JE{k}{t} = 0$ when $\dS{k}{t}$ is undefined.
\end{definition}

Intuitively, $\JE{k}{t} = 1$ means cell $k$ at frame $t$ shows a ``sharp
upward step'' in signal activity that is in the top $(1-q)$ fraction of all
positive changes observed across the block.

% ─────────────────────────────────────────────────────────────────────────────
\section{Spatiotemporal Exposure and Future Response}
% ─────────────────────────────────────────────────────────────────────────────

\subsection{Future self-jump indicator}

\begin{definition}[Future self-jump]\label{def:future-jump}
For node $(k,t)$ and a \emph{future window} $W \in \mathbb{Z}_{>0}$ (frames),
the \emph{future self-jump indicator} is
\[
  \FJ{k}{t} \;=\;
  \1{\exists\, t' \in T_k \;:\; t < t' \leq t + W \;\text{and}\; \JE{k}{t'} = 1}.
\]
Thus $\FJ{k}{t} = 1$ if and only if cell $k$ itself undergoes at least one
jump within the next $W$ frames. The default is $W = 3$, corresponding to a
$15$-minute look-ahead window.
\end{definition}

The future self-jump is computed by a forward scan along each track:
for every node $(k,t)$, the algorithm inspects subsequent frames
$t+1, t+2, \ldots$ until either a jump is found or the window $W$ is
exhausted.

\subsection{Spatial exposure}

\begin{definition}[Neighbourhood statistics and exposure]\label{def:exposure}
For node $(k,t)$, define:
\begin{align}
  n(k,t)       &= |\Nb{k}{t}|
    && \text{(spatial neighbour count),} \label{eq:n}\\[3pt]
  m(k,t)       &= \sum_{(k',t)\,\in\,\Nb{k}{t}} \JE{k'}{t}
    && \text{(number of jumping neighbours at frame $t$),} \label{eq:m}\\[3pt]
  \bar{S}(k,t) &= \frac{1}{n(k,t)}
                   \sum_{(k',t)\,\in\,\Nb{k}{t}} \sig{k'}{t}
    && \text{(mean neighbour signal; undefined if $n(k,t)=0$).} \label{eq:Sbar}
\end{align}
The \emph{exposure indicator} is
\[
  \expo{k}{t} \;=\; \1{m(k,t) > 0},
\]
equal to $1$ when at least one spatial neighbour of cell $k$ is currently
undergoing a jump.
\end{definition}

\begin{remark}
Because spatial edges are rebuilt per frame, $\Nb{k}{t}$ and hence $n(k,t)$
can differ at each time point as cells migrate within the imaging field.
\end{remark}

% ─────────────────────────────────────────────────────────────────────────────
\section{Propagation Statistics}
% ─────────────────────────────────────────────────────────────────────────────

The central question is whether exposure to a jumping neighbour
\emph{right now} is predictive of a future self-jump. We operationalise this
with two classical epidemiological estimators applied to cell-time nodes.

\subsection{Exposed and unexposed populations}

Partition $V$ into two disjoint subsets:
\[
  V_{\mathrm{exp}}   = \{(k,t) \in V \mid \expo{k}{t} = 1\},\qquad
  V_{\mathrm{unexp}} = \{(k,t) \in V \mid \expo{k}{t} = 0\}.
\]

\subsection{Conditional future-jump rates}

Define the \emph{conditional future-jump rates}
\begin{align}
  \pE &= \Pr\!\bigl(\FJ{k}{t}=1 \;\big|\; \expo{k}{t}=1\bigr)
       \;\approx\; \frac{1}{|V_{\mathrm{exp}}|}
         \sum_{(k,t)\,\in\, V_{\mathrm{exp}}} \FJ{k}{t},
       \label{eq:pE}\\[6pt]
  \pU &= \Pr\!\bigl(\FJ{k}{t}=1 \;\big|\; \expo{k}{t}=0\bigr)
       \;\approx\; \frac{1}{|V_{\mathrm{unexp}}|}
         \sum_{(k,t)\,\in\, V_{\mathrm{unexp}}} \FJ{k}{t}.
       \label{eq:pU}
\end{align}

\subsection{Risk difference and relative risk}

\begin{definition}[Risk difference and relative risk]\label{def:RD-RR}
\begin{align}
  \RD &\;=\; \pE \;-\; \pU, \label{eq:RD}\\[4pt]
  \RR &\;=\; \frac{\pE}{\pU} \qquad (\pU > 0). \label{eq:RR}
\end{align}
\end{definition}

\paragraph{Interpretation.}
\begin{itemize}
  \item $\RD = 0$ (or $\RR = 1$): neighbour jumping has no association
        with future self-jumping.
  \item $\RD > 0$ (or $\RR > 1$): exposed cell-time points are more likely
        to be followed by a self-jump.
        An $\RD$ of $0.05$ is a $5$ percentage-point absolute increase.
        An $\RR$ of $1.76$ means exposed cells are $\approx 1.76\times$ more
        likely to jump in the near future.
  \item $\RD < 0$ (or $\RR < 1$): exposure is associated with a
        \emph{lower} future-jump probability (potential suppression).
\end{itemize}

\subsection{Empirical example}

For the WT condition in this dataset (one block, ERK signal):

\begin{center}
\renewcommand{\arraystretch}{1.25}
\begin{tabular}{lrrrrrr}
\toprule
Mutation & $|V|$ & $|E_{\mathrm{sp}}|$ & $\theta$ & $\pE$ & $\pU$ & $\RR$ \\
\midrule
WT & $356{,}632$ & $2{,}389{,}504$ & 0.034 & 0.131 & 0.075 & 1.756 \\
\bottomrule
\end{tabular}
\end{center}

Wild-type cells whose neighbour jumps \emph{now} are roughly $75\%$ more likely
to show their own ERK jump within the next 15 minutes.

\begin{remark}[Statistical note]
Equations \eqref{eq:pE} and \eqref{eq:pU} treat every node as an
independent observation. This is a simplification: nodes from the same track
are serially correlated, and spatially adjacent nodes share neighbourhood
structure. A rigorous inference (e.g.\ permutation test with tracks as units,
or a mixed-effects log-binomial regression) is needed to assign valid
$p$-values to $\RR$.
\end{remark}

% ─────────────────────────────────────────────────────────────────────────────
\section{Multi-Block Comparison Framework}
% ─────────────────────────────────────────────────────────────────────────────

\subsection{Block-level summaries}

An \emph{experiment-site block} $\mathcal{B}_{e,s}$ is the subset of all
observations belonging to experiment $e$ and imaging site $s$.
Let
$\mathcal{B} = \{\mathcal{B}_{e_1,s_1}, \ldots, \mathcal{B}_{e_B,s_B}\}$
be the collection of blocks remaining after filtering (e.g.\ removing the
ErbB2 condition from the comparison).

Running the single-block algorithm of Section~5 on each block $b$ yields the
summary tuple
\[
  \phi(b) \;=\;
  \bigl(n_b,\; |E_{\mathrm{sp},b}|,\; |E_{\mathrm{tm},b}|,\;
    \theta_b,\; \pE^{(b)},\; \pU^{(b)},\; \RD^{(b)},\; \RR^{(b)}\bigr),
\]
where $n_b = |V_b|$ and $\theta_b$ is computed per-block when no global
threshold is supplied.

\subsection{Group-level aggregation}

Let $\gamma : \mathcal{B} \to \Gamma$ assign each block to a comparison
group $g \in \Gamma$ (default: $\Gamma =$ set of distinct mutation labels).
For each group $g$, let $\mathcal{B}_g = \{b \in \mathcal{B} \mid \gamma(b) = g\}$ and define
\begin{align}
  \overline{\RR}(g)
    &= \frac{1}{|\mathcal{B}_g|}
       \sum_{b\,\in\,\mathcal{B}_g} \RR^{(b)},
    \label{eq:meanRR}\\[4pt]
  \widetilde{\RR}(g)
    &= \operatorname{median}_{b\,\in\,\mathcal{B}_g}\bigl\{\RR^{(b)}\bigr\},
    \label{eq:medianRR}\\[4pt]
  \overline{\RD}(g)
    &= \frac{1}{|\mathcal{B}_g|}
       \sum_{b\,\in\,\mathcal{B}_g} \RD^{(b)}.
    \label{eq:meanRD}
\end{align}

Groups are sorted by $\overline{\RR}(g)$ in descending order.
The full aggregated table is saved to \texttt{group\_level\_summary.csv}.

\begin{remark}[Automatic grouping]
When the metadata column \texttt{Conditions} takes only a single value across
all sites (as in this dataset, where all wells share the same growth-factor
deprived condition), the script automatically uses \texttt{Mutation} as the
comparison grouping.
\end{remark}

% ─────────────────────────────────────────────────────────────────────────────
\section{Formal Summary: The Algorithm}
% ─────────────────────────────────────────────────────────────────────────────

\begin{snugshade}
\noindent\textbf{Algorithm~1.} Single-block spatiotemporal propagation analysis.

\medskip
\noindent\textbf{Input:} block $\mathcal{B}_{e,s}$; spatial radius $r$;
  quantile $q$; future window $W$; optional fixed threshold $\theta$.

\noindent\textbf{Output:} summary tuple $\phi(\mathcal{B}_{e,s})$.

\medskip
\begin{enumerate}
  \item \textit{(Signal delta)} For each $(k,t) \in V$:
    compute $\dS{k}{t} = \sig{k}{t} - \sig{k}{t_{\mathrm{prev}}}$
    where $t_{\mathrm{prev}} = \max\{s \in T_k : s < t\}$; set to
    \texttt{NaN} if no predecessor exists.

  \item \textit{(Threshold)} If no fixed $\theta$:
    set $\theta \leftarrow Q_q\!\bigl(\{\dS{k}{t} > 0\}\bigr)$.

  \item \textit{(Jump events)} $\JE{k}{t} \leftarrow \1{\dS{k}{t} \geq \theta}$
    for all $(k,t) \in V$.

  \item \textit{(Spatial edges)} For each frame $t$, build a $k$-d tree on
    $\{\bm{x}(k,t)\}_{t \in T_k}$ and enumerate all pairs within radius $r$
    to form $E_{\mathrm{sp},t}$; set $E_{\mathrm{sp}} = \bigcup_t E_{\mathrm{sp},t}$.

  \item \textit{(Temporal edges)} For each track $k$, connect consecutive
    observations to form $E_{\mathrm{tm}}$.

  \item \textit{(Future self-jump)} For each $(k,t) \in V$:
    scan frames $t+1, t+2, \ldots, t+W$ in $T_k$ and set
    $\FJ{k}{t} \leftarrow 1$ if any of them is a jump event.

  \item \textit{(Exposure)} For each $(k,t) \in V$:
    compute $m(k,t)$ via $E_{\mathrm{sp}}$;
    set $\expo{k}{t} \leftarrow \1{m(k,t)>0}$.

  \item \textit{(Summary)} Compute $\pE$, $\pU$ via
    Eqs.~\eqref{eq:pE}--\eqref{eq:pU};
    return $\RD = \pE - \pU$ and $\RR = \pE/\pU$.
\end{enumerate}
\end{snugshade}

% ─────────────────────────────────────────────────────────────────────────────
\section{Discussion and Connections}
% ─────────────────────────────────────────────────────────────────────────────

\subsection{What this analysis captures}

The analysis measures the \emph{empirical association} between a current
neighbourhood jump and a future self-jump, summarised by $\RR$ and $\RD$.
An $\RR > 1$ is consistent with---but does not prove---mechanisms such as:
\begin{itemize}
  \item direct paracrine signalling (e.g.\ secreted growth factors) between
        physically adjacent cells;
  \item a common upstream spatial stimulus (e.g.\ a localised wave of
        extracellular signal) that activates nearby cells in sequence;
  \item spatially correlated noise in the shared microenvironment.
\end{itemize}

\subsection{Key assumptions and caveats}

\begin{description}
  \item[(a) Node independence.]
        Treating each $(k,t)$ as i.i.d.\ inflates the effective sample size
        and invalidates standard confidence intervals. A valid uncertainty
        estimate uses the \emph{track} as the primary sampling unit.

  \item[(b) Threshold sensitivity.]
        The binary jump indicator depends strongly on $\theta$.
        The per-block quantile threshold makes blocks non-comparable unless
        a global $\theta$ is fixed.

  \item[(c) Temporal resolution.]
        The exposure variable $\expo{k}{t}$ captures \emph{concurrent}
        neighbour jumps only.  A lagged-exposure analysis
        (e.g.\ $\expo{k}{t-1}$ or $\expo{k}{t-2}$) would separate
        temporal precedence from simultaneity.

  \item[(d) Spatial confounding.]
        Spatially adjacent cells may share microenvironmental conditions
        (local ECM density, medium gradients) that independently affect
        signalling, independent of direct cell--cell communication.
\end{description}

\subsection{Connections to other mathematical frameworks}

\begin{description}
  \item[Temporal/dynamic networks.]
        $G = (V, E_{\mathrm{sp}} \cup E_{\mathrm{tm}})$ is a
        \emph{dynamic graph} in which spatial edge sets change at each
        time step while temporal edges encode causal ordering within tracks.

  \item[Geometric random graphs.]
        The spatial edge set $E_{\mathrm{sp},t}$ at any frame $t$ is
        precisely the edge set of the \emph{unit-disc graph} on centres
        $\{\bm{x}(k,t)\}$, i.e.\ the intersection graph of discs of
        radius $r/2$.

  \item[Log-binomial regression.]
        The $\RR$ estimator can be extended to a fully adjusted model:
        \[
          \log \Pr\!\bigl(\FJ{k}{t}=1\bigr)
          = \alpha + \beta\,\expo{k}{t}
            + \bm{\gamma}^\top \bm{z}(k,t),
        \]
        where $\bm{z}(k,t)$ contains covariates such as nuclear size,
        baseline signal, or frame number. Then $e^{\hat\beta} \approx \RR$
        conditionally on those covariates.

  \item[Point process perspective.]
        The sequence of jump events along a track
        $\bigl\{t \in T_k : \JE{k}{t}=1\bigr\}$ is a discrete \emph{point
        process} on the time axis; the analysis asks whether the rate of
        this process is elevated following a jump in a spatially proximal
        process.
\end{description}

% ─────────────────────────────────────────────────────────────────────────────
\section*{Appendix: Notation Summary}
% ─────────────────────────────────────────────────────────────────────────────
\addcontentsline{toc}{section}{Appendix: Notation Summary}

\begin{center}
\renewcommand{\arraystretch}{1.35}
\begin{tabular}{>{\bfseries}p{3.5cm}p{9.5cm}}
\toprule
Symbol & Meaning \\
\midrule
$V$ & Set of all (cell, frame) nodes \\
$E_{\mathrm{sp}}$ & Spatial edges (same frame, within radius $r$) \\
$E_{\mathrm{tm}}$ & Temporal edges (same track, consecutive frames) \\
$k$ & Track (cell) index \\
$t$ & Frame index (integer, $1 \leq t \leq T_{\max}$) \\
$T_k$ & Set of frames at which track $k$ is observed \\
$\bm{x}(k,t)$ & Nuclear centroid $(x,y)$ in image pixels \\
$\sig{k}{t}$ & Signal ratio (ERKKTR or FoxO3A) at node $(k,t)$ \\
$\dS{k}{t}$ & Signal increment: $\sig{k}{t} - \sig{k}{t_{\mathrm{prev}}}$ \\
$\theta$ & Jump threshold (fixed or data-driven) \\
$q$ & Quantile for data-driven $\theta$ (default: $0.90$) \\
$\JE{k}{t}$ & Jump event: $\mathbf{1}[\dS{k}{t} \geq \theta]$ \\
$r$ & Spatial radius for neighbourhood (default: 60 units) \\
$\Nb{k}{t}$ & Spatial neighbourhood of node $(k,t)$ \\
$n(k,t)$ & Neighbour count: $|\Nb{k}{t}|$ \\
$m(k,t)$ & Jumping-neighbour count at frame $t$ \\
$\bar{S}(k,t)$ & Mean signal of spatial neighbours \\
$\expo{k}{t}$ & Exposure indicator: $\mathbf{1}[m(k,t)>0]$ \\
$W$ & Future window in frames (default: 3, i.e.\ 15 min) \\
$\FJ{k}{t}$ & Future self-jump: $\mathbf{1}[\exists\,t'\in(t,t{+}W]:\JE{k}{t'}=1]$ \\
$\pE,\,\pU$ & Conditional future-jump rates (exposed / unexposed) \\
$\RD$ & Risk difference: $\pE - \pU$ \\
$\RR$ & Relative risk: $\pE / \pU$ \\
$\mathcal{B}_{e,s}$ & Experiment-site block (experiment $e$, site $s$) \\
$\overline{\RR}(g)$ & Mean relative risk aggregated over group $g$ \\
\bottomrule
\end{tabular}
\end{center}

\end{document}
